\labsection{1}{R setup, sequences, rounding, and user input}{sec:lab01}

\begin{labproject}{Lab 1b guided practice}
\textbf{Coverage.} Chapter~\ref{chap:r-basics}.\\
\textbf{Goal.} Apply basic sequence generation, arithmetic, rounding, and keyboard input.

\begin{labmode}
Complete the three tasks using only the functions introduced in Chapter~\ref{chap:r-basics}.
\end{labmode}

\textbf{Task 1 --- Create 20 numbers and display their squares.}

\begin{lstlisting}[style=dsr]
numbers <- 1:20
squares <- numbers ^ 2

print(numbers)
print(squares)
\end{lstlisting}

A compact two-column display is:

\begin{lstlisting}[style=dsr]
print(cbind(numbers, squares))
\end{lstlisting}

\textbf{Task 2 --- Round two given values.}

\begin{lstlisting}[style=dsr]
num1 <- 0.956786
num2 <- 7.8345901

print(round(num1, 2))
print(round(num2, 3))
\end{lstlisting}

The expected rounded values are \texttt{0.96} and \texttt{7.835}.

\textbf{Task 3 --- Read a circle radius and calculate area.}

\begin{lstlisting}[style=dsr]
radius <- readline(prompt = "Enter radius: ")
radius <- as.numeric(radius)

area <- pi * radius ^ 2
print(paste("Circle area:", area))
\end{lstlisting}

The program reads, converts, calculates, and displays.
\end{labproject}

\begin{labdeliverable}
Submit all three programs and verified outputs to your instructor.
\end{labdeliverable}
